% from ega1-4-fr.tex
\oldpage[I]{126}
de façon plus générale dans (\hyperref[I.3.3.6-fr]{3.3.6}). Quand nous
parlerons de $f^{-1}(Y')$ comme d'un sous-préschéma de $X$, c'est toujours de
ce sous-préschéma qu'il s'agira.

Lorsque les préschémas $f^{-1}(Y')$ et $X$ sont identiques, $j'$ est
l'identité et tout morphisme $h:Z\to X$ est donc majoré par $j'$, donc
\emph{le morphisme $f:X\to Y$ se factorise alors en}
$X\xrightarrow{g}Y'\xrightarrow{j}Y$.

Lorsque $y$ est un point \emph{fermé} de $Y$ et
$Y'=\operatorname{Spec}(k(y))$ le plus petit sous-préschéma fermé de $Y$ ayant
$\{y\}$ comme espace sous-jacent (\hyperref[I.4.1.9-fr]{4.1.9}), le
sous-préschéma fermé $f^{-1}(Y')$ est canoniquement isomorphe à la
\emph{fibre} $f^{-1}(y)$ définie en (\hyperref[I.3.6.2-fr]{3.6.2}), à
laquelle on l'identifiera.

\begin{corollary}[4.4.2]
\label{I.4.4.2-fr}
Soient $f:X\to Y$, $g:Y\to Z$ deux morphismes, $h=g\circ f$ leur composé.
Pour tout sous-préschéma $Z'$ de $Z$, les sous-préschémas
$f^{-1}(g^{-1}(Z'))$ et $h^{-1}(Z')$ de $X$ sont identiques.
\end{corollary}

Cela résulte de l'existence de l'isomorphisme canonique
$X\times_Y(Y\times_Z Z')\xrightarrow{\sim}X\times_Z Z'$
(\hyperref[I.3.3.9.1-fr]{3.3.9.1}).

\begin{corollary}[4.4.3]
\label{I.4.4.3-fr}
Soient $X'$, $X''$ deux sous-préschémas de $X$, $j':X'\to X$,
$j'':X''\to X$ les morphismes d'injection~; alors, ${j'}^{-1}(X'')$ et
${j''}^{-1}(X')$ sont tous deux égaux à la borne inférieure
$\inf(X',X'')$ de $X'$ et $X''$ pour la relation d'ordre entre
sous-préschémas, et canoniquement isomorphée à $X'\times_X X''$.
\end{corollary}

Cela résulte aussitôt de (\hyperref[I.4.4.1-fr]{4.4.1}) et
(\hyperref[I.4.1.10-fr]{4.1.10}).

\begin{corollary}[4.4.4]
\label{I.4.4.4-fr}
Soient $f:X\to Y$ un morphisme, $Y'$, $Y''$ deux sous-préschémas de $Y$~;
on a
\[
  f^{-1}(\inf(Y',Y''))=\inf(f^{-1}(Y'),f^{-1}(Y'')).
\]
\end{corollary}

Cela résulte de l'existence de l'isomorphisme canonique entre
$(X\times_Y Y')\times_X(X\times_Y Y'')$ et
$X\times_Y(Y'\times_Y Y'')$ (\hyperref[I.3.3.9.1-fr]{3.3.9.1}).

\begin{proposition}[4.4.5]
\label{I.4.4.5-fr}
Soient $f:X\to Y$ un morphisme, $Y'$ un sous-préschéma fermé de $Y$ défini
par un faisceau quasi-cohérent $\mathscr{K}$ d'idéaux de $\mathscr{O}_Y$
(\hyperref[I.4.1.3-fr]{4.1.3})~; le sous-préschéma fermé $f^{-1}(Y')$ de $X$
est alors défini par le faisceau quasi cohérent d'idéaux
$f^*(\mathscr{K})\mathscr{O}_X$ de $\mathscr{O}_X$.
\end{proposition}

La question est évidemment locale sur $X$ et $Y$~; il suffit alors de
remarquer que si $B$ est une $A$-algèbre et $\mathfrak{K}$ un idéal de $B$,
on a $A\otimes_B(B/\mathfrak{K})=A/\mathfrak{K}A$, et d'appliquer
(\hyperref[I.1.6.9-fr]{1.6.9}).

\begin{corollary}[4.4.6]
\label{I.4.4.6-fr}
Soient $X'$ un sous-préschéma fermé de $X$ défini par un faisceau quasi
cohérent d'idéaux $\mathscr{J}$ de $\mathscr{O}_X$, $i$ l'injection
$X'\to X$~; pour que la restriction $f\circ i$ de $f$ à $X'$ soit majorée
par l'injection $j:Y'\to Y$ (autrement dit se factorise en $j\circ g$, où
$g$ est un morphisme $X'\to Y'$), il faut et il suffit que
$f^*(\mathscr{K})\mathscr{O}_X\subset\mathscr{J}$.
\end{corollary}

Il suffit d'appliquer à $i$ la prop. (\hyperref[I.4.4.1-fr]{4.4.1}) en
tenant compte de (\hyperref[I.4.4.5-fr]{4.4.5}).

\subsection{Immersions locales et isomorphismes locaux.}
\label{subsection:I.4.5-fr}

\begin{definition}[4.5.1]
\label{I.4.5.1-fr}
Soit $f:X\to Y$ un morphisme de préschémas. On dit que $f$ est une
\emph{immersion locale en un point $x\in X$} s'il existe un voisinage ouvert
$U$ de $x$ dans $X$ et un voisinage ouvert $V$ de $f(x)$ dans $Y$ tels que
la restriction de $f$ au préschéma induit $U$ soit une immersion fermée de
$U$ dans le préschéma induit $V$. On dit que $f$ est une immersion locale si
$f$ est une immersion locale en tout point de $X$.
\end{definition}

\begin{definition}[4.5.2]
\label{I.4.5.2-fr}
On dit qu'un morphisme $f:X\to Y$ est un isomorphisme local en
